\documentclass{amsart}
% For dealing with references we use the comment environment
\usepackage{verbatim}
\newenvironment{reference}{\comment}{\endcomment}
%\newenvironment{reference}{}{}
\newenvironment{slogan}{\comment}{\endcomment}
\newenvironment{history}{\comment}{\endcomment}

% For commutative diagrams we use Xy-pic
\usepackage[all]{xy}

% We use 2cell for 2-commutative diagrams.
\xyoption{2cell}
\UseAllTwocells

% We use multicol for the list of chapters between chapters
\usepackage{multicol}

% This is generally recommended for better output
\usepackage{lmodern}
\usepackage[T1]{fontenc}

% For cross-file-references

% Package for hypertext links:
\usepackage{hyperref}

% For any local file, say "hello.tex" you want to link to please

% Theorem environments.
%
\theoremstyle{plain}
\newtheorem{theorem}[subsection]{Theorem}
\newtheorem{proposition}[subsection]{Proposition}
\newtheorem{lemma}[subsection]{Lemma}

\theoremstyle{definition}
\newtheorem{definition}[subsection]{Definition}
\newtheorem{example}[subsection]{Example}
\newtheorem{exercise}[subsection]{Exercise}
\newtheorem{situation}[subsection]{Situation}

\theoremstyle{remark}
\newtheorem{remark}[subsection]{Remark}
\newtheorem{remarks}[subsection]{Remarks}

\numberwithin{equation}{subsection}

% Macros
%
\def\lim{\mathop{\mathrm{lim}}\nolimits}
\def\colim{\mathop{\mathrm{colim}}\nolimits}
\def\Spec{\mathop{\mathrm{Spec}}}
\def\Hom{\mathop{\mathrm{Hom}}\nolimits}
\def\Ext{\mathop{\mathrm{Ext}}\nolimits}
\def\SheafHom{\mathop{\mathcal{H}\!\mathit{om}}\nolimits}
\def\SheafExt{\mathop{\mathcal{E}\!\mathit{xt}}\nolimits}
\def\Sch{\mathit{Sch}}
\def\Mor{\mathop{\mathrm{Mor}}\nolimits}
\def\Ob{\mathop{\mathrm{Ob}}\nolimits}
\def\Sh{\mathop{\mathit{Sh}}\nolimits}
\def\NL{\mathop{N\!L}\nolimits}
\def\CH{\mathop{\mathrm{CH}}\nolimits}
\def\proetale{{pro\text{-}\acute{e}tale}}
\def\etale{{\acute{e}tale}}
\def\QCoh{\mathit{QCoh}}
\def\Ker{\mathop{\mathrm{Ker}}}
\def\Im{\mathop{\mathrm{Im}}}
\def\Coker{\mathop{\mathrm{Coker}}}
\def\Coim{\mathop{\mathrm{Coim}}}

% Boxtimes
%
\DeclareMathSymbol{\boxtimes}{\mathbin}{AMSa}{"02}

%
% Macros for moduli stacks/spaces
%
\def\QCohstack{\mathcal{QC}\!\mathit{oh}}
\def\Cohstack{\mathcal{C}\!\mathit{oh}}
\def\Spacesstack{\mathcal{S}\!\mathit{paces}}
\def\Quotfunctor{\mathrm{Quot}}
\def\Hilbfunctor{\mathrm{Hilb}}
\def\Curvesstack{\mathcal{C}\!\mathit{urves}}
\def\Polarizedstack{\mathcal{P}\!\mathit{olarized}}
\def\Complexesstack{\mathcal{C}\!\mathit{omplexes}}
% \Pic is the operator that assigns to X its picard group, usage \Pic(X)
% \Picardstack_{X/B} denotes the Picard stack of X over B
% \Picardfunctor_{X/B} denotes the Picard functor of X over B
\def\Pic{\mathop{\mathrm{Pic}}\nolimits}
\def\Picardstack{\mathcal{P}\!\mathit{ic}}
\def\Picardfunctor{\mathrm{Pic}}
\def\Deformationcategory{\mathcal{D}\!\mathit{ef}}

\setlength{\emergencystretch}{3em}
\begin{document}
\title[Stacks Project, Tag 0ASI]{425 --- Descent of roots up to units through codimension-one unramified maps}
\maketitle
\noindent Copyright (C) 2005--2025 Johan de Jong. Stacks Project authors. Source: \href{https://stacks.math.columbia.edu/tag/0ASI}{Tag 0ASI}.

\noindent Editorial difficulty note (not author text). Difficulty H2: The proof constructs a fractional-valuation ideal from the required divided orders at height-one primes. Tensoring and normal-domain intersection criteria show that ideal becomes principal upstairs, and faithfully flat descent supplies the desired root downstairs..

\noindent Editorial provenance note (not author text). History: excerpt from revision \texttt{a0444\allowbreak 6e57e\allowbreak c1fbc\allowbreak 252a8\allowbreak 71afc\allowbreak ec775\allowbreak 2fb28\allowbreak 07b14}, more-algebra.tex, lines 37390--37477. Reference presentation was adapted; original-author context and core support blocks were retained or added. The mathematical wording is unchanged. Licensed under GNU FDL 1.2 or later; no invariant sections or cover texts. See ../licenses/COPYING. Context blocks reproduce author definitions and supporting results; they are not additional counted problems.

\begin{definition}
\label{definition-extension-discrete-valuation-rings}
We say that $A \to B$ or $A \subset B$ is an
{\it extension of discrete valuation rings} if $A$ and $B$ are
discrete valuation rings and $A \to B$ is injective and local.
In particular, if $\pi_A$ and $\pi_B$ are uniformizers of
$A$ and $B$, then $\pi_A = u \pi_B^e$ for some $e \geq 1$ and unit
$u$ of $B$. The integer $e$ does not depend on the choice of
the uniformizers as it is also the unique integer $\geq 1$ such that
$$
\mathfrak m_A B = \mathfrak m_B^e
$$
The integer $e$ is called the {\it ramification index} of $B$ over $A$.
We say that $B$ is {\it weakly unramified} over $A$ if $e = 1$.
If the extension of residue fields
$\kappa_A = A/\mathfrak m_A \subset \kappa_B = B/\mathfrak m_B$
is finite, then we set $f = [\kappa_B : \kappa_A]$ and we
call it the {\it residual degree} or {\it residue degree}
of the extension $A \subset B$.
\end{definition}

\begin{lemma}
\label{lemma-extension-normal-domains-and-roots}
Let $A \to B$ be a flat local homomorphism of Noetherian local normal domains.
Let $f \in A$ and $h \in B$ such that $f = w h^n$ for some $n > 1$ and some
unit $w$ of $B$. Assume that for every height $1$ prime
$\mathfrak p \subset A$ there is a height $1$ prime
$\mathfrak q \subset B$ lying over $\mathfrak p$
such that the extension $A_\mathfrak p \subset B_\mathfrak q$ is
weakly unramified. Then $f = u g^n$ for some $g \in A$ and unit $u$ of $A$.
\end{lemma}

\begin{proof}
The local rings of $A$ and $B$ at height $1$ primes are
discrete valuation rings (Algebra, Lemma \href{https://stacks.math.columbia.edu/tag/00PD}{lemma characterize dvr (Tag 00PD)}).
Thus the assumption makes sense (via
Definition \ref{definition-extension-discrete-valuation-rings}).
Let $\mathfrak p_1, \ldots, \mathfrak p_r$ be the primes of $A$ minimal
over $f$. These have height $1$ by
Algebra, Lemma \href{https://stacks.math.columbia.edu/tag/00KV}{lemma minimal over 1 (Tag 00KV)}.
For each $i$ let $\mathfrak q_{i, j} \subset B$, $j = 1, \ldots, r_i$
be the height $1$ primes of $B$ lying over $\mathfrak p_i$.
Say we number them so that  $A_{\mathfrak p_i} \to B_{\mathfrak q_{i, 1}}$
is weakly unramified.
Since $f$ maps to an $n$th power times a unit in $B_{\mathfrak q_{i, 1}}$
we see that the valuation $v_i$ of $f$ in $A_{\mathfrak p_i}$ is
divisible by $n$. Say $v_i = n w_i$ for some $w_i \geq 0$.
Consider the exact sequence
$$
0 \to I \to A \to
\prod\nolimits_{i = 1, \ldots, r}
A_{\mathfrak p_i}/\mathfrak p_i^{w_i}A_{\mathfrak p_i}
$$
defining the ideal $I$. Applying the exact functor $- \otimes_A B$ we obtain
an exact sequence
$$
0 \to I \otimes_A B \to B \to
\prod\nolimits_{i = 1, \ldots, r}
(A_{\mathfrak p_i}/\mathfrak p_i^{w_i}A_{\mathfrak p_i}) \otimes_A B
$$
Fix $i$. We claim that the canonical map
$$
(A_{\mathfrak p_i}/\mathfrak p_i^{w_i}A_{\mathfrak p_i}) \otimes_A B
\to
\prod\nolimits_{j = 1, \ldots, r_i}
B_{\mathfrak q_{i, j}}/\mathfrak q_{i, j}^{e_{i, j}w_i}B_{\mathfrak q_{i, j}}
$$
is injective. Here $e_{i, j}$ is the ramification index of 
$A_{\mathfrak p_i} \to B_{\mathfrak q_{i, j}}$.
The claim asserts that $\mathfrak p_i^{w_i}B_{\mathfrak p_i}$
is equal to the set of elements $b$ of $B_{\mathfrak p_i}$
whose valuation at $\mathfrak q_{i, j}$ is $\geq e_{i, j}w_i$.
Choose a generator $a \in A_{\mathfrak p_i}$ of the principal
ideal $\mathfrak p_i^{w_i}$. Then the valuation of $a$ at
$\mathfrak q_{i, j}$ is equal to $e_{i, j}w_i$. Hence,
as $B_{\mathfrak p_i}$ is a normal domain whose height one
primes are the primes $\mathfrak q_{i, j}$, $j = 1, \ldots, r_i$,
we see that, for $b$ as above, we have $b/a \in B_{\mathfrak p_i}$ by
Algebra, Lemma
\href{https://stacks.math.columbia.edu/tag/031T}{lemma normal domain intersection localizations height 1 (Tag 031T)}.
Thus the claim.

\medskip\noindent
The claim combined with the second exact sequence above
determines an exact sequence
$$
0 \to I \otimes_A B \to B \to
\prod\nolimits_{i = 1, \ldots, r}
\prod\nolimits_{j = 1, \ldots, r_i}
B_{\mathfrak q_{i, j}}/\mathfrak q_{i, j}^{e_{i, j}w_i}B_{\mathfrak q_{i, j}}
$$
It follows that $I \otimes_A B$ is the set of elements $h'$ of $B$
which have valuation $\geq e_{i, j}w_i$ at $\mathfrak q_{i, j}$.
Since $f = wh^n$ in $B$ we see that $h$ has valuation
$e_{i, j}w_i$ at $\mathfrak q_{i, j}$. Thus $h'/h \in B$
by Algebra, Lemma
\href{https://stacks.math.columbia.edu/tag/031T}{lemma normal domain intersection localizations height 1 (Tag 031T)}.
It follows that $I \otimes_A B$ is a free $B$-module of rank $1$
(generated by $h$).
Therefore $I$ is a free $A$-module of rank $1$, see
Algebra, Lemma \href{https://stacks.math.columbia.edu/tag/00O1}{lemma finite projective descends (Tag 00O1)}.
Let $g \in I$ be a generator. Then we see that
$g$ and $h$ differ by a unit in $B$. Working backwards we
conclude that the valuation of $g$ in $A_{\mathfrak p_i}$ is
$w_i = v_i/n$. Hence $g^n$ and $f$ differ by a unit in $A$
(by Algebra, Lemma
\href{https://stacks.math.columbia.edu/tag/031T}{lemma normal domain intersection localizations height 1 (Tag 031T)})
as desired.
\end{proof}
\end{document}
